\documentclass[journal]{IEEEtran}
\usepackage{amsmath,amsfonts}
\usepackage{algorithmic}
\usepackage{algorithm}
\usepackage{array}
\usepackage[caption=false,font=normalsize,labelfont=sf,textfont=sf]{subfig}
\usepackage{textcomp}
\usepackage{stfloats}
\usepackage{url}
\usepackage{verbatim}
\usepackage{graphicx}
\usepackage{cite}
\usepackage[hidelinks]{hyperref}
\newcommand{\artifact}[1]{\path{#1}}
\hyphenation{op-tical net-works semi-conduc-tor IEEE-Xplore}

\begin{document}

\title{Adaptive Risk-Aware and Explainable Typical-Day Library Learning for Transferable Renewable Dispatch: An IEEE-33 Learning Benchmark With Mapped-System AC-OPF Validation}

\author{Zhaoyi Sun, Ping Liu, and Peipei Meng%
\thanks{Zhaoyi Sun is with EFREI Paris, Panth\'eon-Assas University, 94800 Villejuif, France, and also with the School of Information Engineering, Wuhan University of Technology, 430070 Wuhan, China.}%
\thanks{Ping Liu is with the School of Information Engineering, Wuhan University of Technology, 430070 Wuhan, China.}%
\thanks{Peipei Meng is with the School of Automation, Wuhan University of Technology, Wuhan 430070, China, and is the corresponding author (e-mail: cathleen@whut.edu.cn).}}

\markboth{IEEE Transactions on Smart Grid,~under review}%
{Sun2Paper: Adaptive Risk-Aware and Explainable Typical-Day Library Learning}

\maketitle

\begin{abstract}
We introduce an adaptive risk-aware and explainable representative-day library framework for renewable dispatch, where reduced scenario sets are evaluated by operational stress rather than by geometric similarity alone. The framework combines a learnable risk model, weighted xRFM embedding, dual-layer representative selection, SHAP-based consistency checking, replay-stabilized updating, and two complementary physical-validation layers. It is designed as a dispatch-facing scenario-management front end rather than as a solver replacement. On DDRE-33, the proposed dual-layer xRFM attains the lowest five-seed transfer proxy error, $0.3925 \pm 0.0229$, with a 95\% confidence interval of [0.3641, 0.4210], compared with $0.4593 \pm 0.0281$ for hierarchical clustering, $0.4646 \pm 0.0247$ for spectral clustering, $0.4739 \pm 0.0215$ for single-layer xRFM, $0.4787 \pm 0.0323$ for K-means, and $0.4806 \pm 0.0353$ for k-medoids. The same library also achieves the smallest feeder-aware dispatch-cost gap on IEEE-33 (0.0607) and on an IEEE-69-like scaled feeder (0.0620). To test whether the learned library preserves native feeder stress, we add a DDRE-33 radial branch-flow screen on the 33-bus topology. Under this direct check, the proposed library yields the lowest daily minimum-voltage MAE, $1.15\times10^{-3}$ p.u., and the lowest daily-loss MAE, 0.098 MWh, among the tested methods. We further formulate a full AC-OPF benchmark with explicit nodal power-balance equations, branch-loading constraints, generator limits, and voltage bounds, and evaluate reduced scenario sets on IEEE 30-bus and 118-bus standard transmission systems using DDRE-derived net-load scenarios. All AC-OPF runs converge, with mean power-balance residuals of $1.57\times10^{-6}$ MW and $2.60\times10^{-5}$ MW for IEEE-30 and IEEE-118, respectively. A 20-draw repeated stratified-sampling OPF-reduction study shows that the best retained sets keep mean objective MAPE near 1.0\% at $K=4$ and below 0.3\% at $K=12$. An added native-feeder snapshot AC-OPF study on DDRE-33 gives 100\% convergence for full and representative scenarios, with the proposed library achieving the lowest line-loading MAE (0.399\%) and a low minimum-voltage MAE ($1.49\times10^{-3}$ p.u.); on heterogeneous IEEE-69, it again gives the lowest line-loading MAE (0.479\%) and minimum-voltage MAE ($2.96\times10^{-4}$ p.u.).
\end{abstract}

\begin{IEEEkeywords}
Representative-day learning, scenario reduction, renewable dispatch, explainable learning, continual updating, AC optimal power flow, smart grid validation.
\end{IEEEkeywords}

\section{Introduction}
\IEEEPARstart{T}{he} rapid growth of distributed wind and photovoltaic resources has made distribution-system operation more volatile, less stationary, and more exposed to rare ramps, prolonged low-output intervals, and cross-node mismatch events. In practice, operators do not dispatch directly on every historical or forecast renewable trajectory. A more realistic workflow is to compress a large scenario set into a representative library, solve a reduced scheduling problem, and update that library when seasonal drift or operating conditions change. Under such a workflow, representative-day learning is valuable only when it preserves dispatch-relevant operating difficulty rather than geometric compactness alone.

Existing representative-period and scenario-compression methods have produced strong results for planning-oriented approximation, clustering compactness, and time-series aggregation [1]--[7]. Yet those objectives do not ensure that the scenarios driving dispatch difficulty are retained. A library can be compact in feature space and still miss the days that matter most to the downstream operating problem. Likewise, explainable learning improves transparency [8], and continual-learning techniques improve temporal adaptation [9]--[11], but these ingredients are seldom combined in one representative-library design aimed at dispatch transfer.

The earlier version of this study focused on scenario representation, risk-aware library construction, and surrogate dispatch validation. Reviewer feedback exposed three weaknesses: the physical validation was too indirect, the comparison with OPF-coupled reduction was too limited, and the claim boundary was not sharp enough. The current version addresses these issues directly. It adds a direct DDRE-33 radial branch-flow screening layer, closes native-feeder snapshot AC-OPF on the Baran-Wu 33-bus test feeder, formulates a full AC-OPF benchmark on IEEE-30 and IEEE-118, extends the OPF-reduction analysis from a single benchmark panel to 20 repeated stratified draws, and adds a heterogeneous-feeder snapshot AC-OPF supplement on a 69-bus feeder. As a result, the paper makes a narrower claim than before, but supports that claim with stronger evidence.

Our contributions are as follows:
\begin{itemize}
\item We introduce a unified representative-day learning framework that jointly adapts risk scoring, weighted xRFM embedding, dual-layer representative selection, SHAP-based consistency checking, and replay-based updating to the dispatch-relevant operating difficulty of renewable scenarios.
\item We add a direct DDRE-33 radial screening layer that tests whether reduced scenario libraries preserve native-feeder voltage, branch-loading, and loss patterns.
\item We formulate a mapped-system AC-OPF benchmark on IEEE-30 and IEEE-118 and strengthen the OPF-reduction evidence with repeated stratified sampling.
\item We keep the claim boundary explicit: the method is a dispatch-facing scenario-management front end, not a full AC distribution OPF replacement on DDRE-33.
\end{itemize}
Together, these elements connect data-driven scenario compression to operational behavior that matters for power-system applications.

One scope boundary is important from the outset. The manuscript remains centered on DDRE-33 representative-day learning. The added DDRE-33 radial screen evaluates the native feeder object directly, and the new 33-bus snapshot AC-OPF layer adds executable native-feeder optimization evidence at the single-snapshot level. The IEEE-30/118 AC-OPF study therefore serves as a bridge from data-driven compression to explicit network constraints under a benchmark setting, while the IEEE-69 snapshot AC-OPF supplement provides heterogeneous-feeder evidence under a distribution setting. Keeping the distinction between native-feeder snapshot closure and broader multiperiod/device-rich distribution OPF explicit keeps the contribution, the evidence, and the remaining limitations aligned.

\section{Related Work}
Representative-period learning for energy systems has developed along several lines that only partly overlap. One line studies geometry-oriented representative days and time-series aggregation [1]--[7]. These methods are useful for compression and approximation, but they mostly reward compactness or reconstruction quality. A second line studies problem-driven or operationally coupled scenario reduction, where scenario importance is defined by its influence on the downstream optimization problem rather than by geometric distance alone [14]. A third line studies direct data-driven OPF or scheduling surrogates, which learn the mapping from uncertain operating states to dispatch decisions more explicitly [15]. A fourth line brings explainability and continual adaptation into power-system learning tasks [8]--[11].

Our work sits between these lines rather than squarely inside any one of them. It is closer to operation than geometry-oriented aggregation because it weights scenarios by dispatch-facing difficulty and tests the selected library under feeder screening and AC-OPF validation. At the same time, it remains less tightly coupled to the downstream solver than problem-driven reduction or direct OPF learning because it is still an upstream scenario-library layer. That distinction matters because the goal here is not to replace the dispatch model, but to produce a representative library that remains reusable, interpretable, and updateable while preserving the operating regimes that downstream dispatch actually sees.

This is also where the novelty should be read. Clustering, SHAP, replay updating, and OPF validation are not individually new. What is new is the task-specific coordination of these components around one concrete objective: preserve dispatch-relevant operating difficulty in a representative-day library, explain the learned emphasis, stabilize the library under drift, and test the result through both native-feeder screening and mapped-system AC-OPF. In other words, the contribution lies less in proposing a new primitive model than in converting several general AI tools into a physically disciplined, dispatch-facing representative-library workflow. That integration gives the paper a different role from either pure aggregation work or full solver-facing learning.

The positioning is intentionally moderate. Problem-driven reduction can be more tightly optimized for a specific downstream objective [14], and direct multiperiod OPF learning can be physically deeper [15]. Our claim is narrower: representative-day libraries still have value when the desired artifact is a transparent scenario front end that can serve more than one downstream planning or dispatch workflow.

\section{Problem Formulation}
\subsection{Risk-Aware Representative-Day Objective}
Each daily integrated renewable scenario is denoted by $x_i(t)\in\mathbb{R}^{4}$ for $t=1,\dots,96$, where the four channels correspond to two wind nodes and two PV nodes in DDRE-33. A feature vector $\phi_i\in\mathbb{R}^d$ summarizes mean output, peak level, volatility, ramp intensity, smoothness, complementarity, and context labels. The learnable risk score is
\begin{equation}
r_i=(w^\top q_i)\Big(1+\gamma\frac{|\mu_i-\mu_{ref}|}{\mu_{ref}+\varepsilon}\Big),
\end{equation}
where $q_i$ contains $\mathrm{RampMax}_i$, $\mathrm{Smoothness}_i$, $\mathrm{LowOutputRate}_i$, and $\mathrm{PowerDeficit}_i$, and $w$ is fitted through transfer-oriented proxy supervision:
\begin{equation}
w^\star=\arg\min_w \sum_{i\in\mathcal{D}_{train}} \|\hat{J}^{proxy}_i(w)-J^{proxy}_i\|_2^2+\rho\|w\|_2^2.
\end{equation}
The representative-learning objective is
\begin{equation}
z_i=W^\top \tilde{\phi}_i,\qquad \tilde{\phi}_i=\frac{\phi_i-\bar{\phi}}{\sigma_\phi},
\end{equation}
\begin{equation}
\begin{aligned}
\min_{W,S,w}\quad & \sum_i (1+\alpha r_i)\|z_i-z_{c(i)}\|_2^2 \\
& +\beta \mathcal{L}_{cov}(S)+\eta \mathcal{L}_{stab}(S)+\rho\|w\|_2^2,
\end{aligned}
\end{equation}
where $S=\mathcal{L}_n\cup\mathcal{L}_r$ is the dual-layer library, $\mathcal{L}_{cov}$ penalizes insufficient coverage, and $\mathcal{L}_{stab}$ penalizes unstable updates under replay.

The purpose of the risk term is methodological rather than cosmetic. Rare ramps, sustained low-output windows, and deficit-heavy days can have small influence on average reconstruction error yet large influence on downstream dispatch difficulty. Weighting the embedding loss by $1+\alpha r_i$ shifts representation pressure toward those dispatch-difficult scenarios. Likewise, the dual-layer structure is not simply a larger library. $\mathcal{L}_n$ is used to preserve dense coverage of regular operating states, whereas $\mathcal{L}_r$ acts as a protected tail set so that infrequent but operationally important days are not absorbed into normal clusters. In the mathematical limit $\alpha\to\infty$, the weighted term by itself would increasingly favor tail scenarios; however, the implemented problem does not permit the library to collapse into an all-extreme set because $\mathcal{L}_n$ is retained explicitly and $\mathcal{L}_{cov}$ continues to penalize loss of normal-state coverage. We therefore interpret large $\alpha$ as a stronger tail-emphasis setting inside a constrained dual-layer covering problem, rather than as a regime in which normal operating modes disappear from the selected library.

\subsection{Physical Validation Layers}
To verify that reduced scenario sets remain meaningful under explicit network constraints, we introduce a full AC-OPF benchmark on standard transmission test systems. For each bus $i\in\mathcal{N}$, the model minimizes total generation cost
\begin{equation}
\min \sum_{g\in\mathcal{G}} C_g(P_g),
\end{equation}
subject to nodal active-power balance
\begin{equation}
\begin{aligned}
P_{Gi}-P_{Di}=V_i\sum_{j\in\mathcal{N}} V_j(&G_{ij}\cos(\theta_i-\theta_j) \\
& +B_{ij}\sin(\theta_i-\theta_j)).
\end{aligned}
\end{equation}
and nodal reactive-power balance
\begin{equation}
\begin{aligned}
Q_{Gi}-Q_{Di}=V_i\sum_{j\in\mathcal{N}} V_j(&G_{ij}\sin(\theta_i-\theta_j) \\
& -B_{ij}\cos(\theta_i-\theta_j)).
\end{aligned}
\end{equation}
Generator limits are enforced by
\begin{equation}
P_{Gi}^{\min}\le P_{Gi}\le P_{Gi}^{\max},\qquad
Q_{Gi}^{\min}\le Q_{Gi}\le Q_{Gi}^{\max},
\end{equation}
voltage bounds by
\begin{equation}
V_i^{\min}\le V_i\le V_i^{\max},
\end{equation}
and branch apparent-power limits by
\begin{equation}
|S_{ij}|\le S_{ij}^{\max},\qquad \forall (i,j)\in\mathcal{E}.
\end{equation}
The benchmark is solved with a MATPOWER/PYPOWER-style AC OPF implementation [18] on IEEE-30 and IEEE-118. For each DDRE-derived scenario, system load is scaled by $\lambda_i^{load}$, reactive demand by $\lambda_i^{reactive}$, and renewable injections are mapped to selected PQ buses as negative net load. This design preserves the statistical structure of DDRE scenarios while making them executable on standard network cases.

To partially close the native-feeder gap, we also add a direct DDRE-33 radial validation layer. Using the standard Baran-Wu 33-bus topology, actual renewable injections at buses 18, 22, 25, and 33, and a daily load-shape multiplier, each scenario is screened through a backward-forward branch-flow sweep. For each day, we record the daily minimum voltage magnitude, the daily maximum branch loading, the daily feeder loss, and the fraction of time steps satisfying the adopted voltage and line-loading bounds. This layer is weaker than full AC distribution OPF because it does not optimize reactive support, tap positions, or dispatch controls, but it directly tests whether representative days preserve feeder-level stress patterns on the DDRE-33 network object itself.

\subsection{OPF-Driven Scenario Reduction Objective}
Given a full scenario set $\Omega$, we seek a reduced set $\Omega_K\subset\Omega$, $|\Omega_K|=K$, that preserves the OPF response of the full set. The validation criterion is not merely feature reconstruction. Instead, reduced-set quality is evaluated through AC-OPF objective error, branch-loading error, voltage-profile error, and computation-time reduction. Two OPF-driven methods are added:
\begin{enumerate}
\item Sensitivity-driven selection, which ranks scenarios by deviation in OPF objective, branch loading, voltage stress, and active-constraint activation.
\item KKT-inspired active-set selection, which groups scenarios by near-binding branch, voltage, and generator constraints and retains representatives from the most influential active-set patterns.
\end{enumerate}
These are compared against K-means representative selection and forward feature-space selection. The second OPF-driven method is intentionally described as KKT-inspired rather than exact-dual-driven. A strict KKT treatment would compare scenarios through continuous dual multipliers and complementarity-consistent active sets. Our implementation relaxes that requirement by replacing exact dual-variable matching with thresholded near-binding branch, voltage, and generator signatures, and then clustering scenarios by the resulting discrete activation-pattern distance. In other words, active-set similarity is approximated through observed constraint-activation patterns rather than exact shadow-price geometry in dual space.

\section{Proposed Method}
\subsection{Unified Learning Pipeline}
The framework contains six coupled stages:
\begin{enumerate}
\item Feature engineering transforms each day-level multivariate sequence into descriptive and risk-related variables.
\item A learnable risk model estimates transfer-oriented scenario importance.
\item Weighted xRFM embedding assigns larger influence to scenarios with higher dispatch relevance.
\item Dual-layer selection builds $\mathcal{L}_n$ for regular operating states and $\mathcal{L}_r$ for high-risk scenarios above threshold $\tau$.
\item SHAP interpretation checks whether high-importance features coincide with dispatch-difficult scenarios.
\item Replay-based updating refreshes the library when weekly batches arrive while reducing forgetting.
\end{enumerate}

\subsection{Validation Layers}
The physical-validation layers sit outside the library-learning objective, but they are central to how the learned library is judged. For each selected or full scenario, the DDRE-33 radial screen or the mapped AC-OPF benchmark records:
\begin{itemize}
\item objective value
\item minimum voltage magnitude
\item maximum branch loading
\item daily feeder loss
\item power-balance residual
\item active-constraint signature
\item solve time
\end{itemize}
This allows reduced scenario sets to be judged by operational fidelity rather than by geometric proximity alone.

Two interpretation boundaries are important. First, SHAP is used as a consistency check on the learned selection logic rather than as an optimization variable inside the library-learning objective. Its attribution values are local, sample-level explanations of why an individual day is judged influential by the learned risk-and-selection model; they are not a proof that the final representative library is globally optimal as a set. The global adequacy of the selected library is instead enforced by the weighted embedding, the clustering-and-coverage objective, and the dual-layer construction $S=\mathcal{L}_n\cup\mathcal{L}_r$. Second, neither the DDRE-33 radial screening layer nor the mapped AC-OPF benchmark is a replacement for the upstream scenario-learning module. This separation is deliberate: it keeps the method reusable across downstream workflows while still testing whether the selected scenarios preserve constraint-sensitive behavior.

\subsection{Algorithmic Complexity}
Feature extraction requires $O(NTd)$, where $N$ is scenario count, $T=96$, and $d$ is feature dimension. Risk-weight fitting contributes $O(Nb^2+b^3)$ with low-dimensional risk basis size $b$. Weighted embedding and clustering contribute approximately $O(Nd^2+d^3+NKdI)$, where $K$ is library size and $I$ is cluster-iteration count. SHAP surrogate fitting contributes $O(BNd\log N)$ for $B$ trees. Replay updating contributes $O(R(M+b)d)$, where $M$ is replay size and $R$ is update rounds.

For the new AC-OPF validation, solving $N$ scenarios on one test system contributes $O(N \cdot T_{OPF})$, where $T_{OPF}$ denotes average AC-OPF solve time. Empirically, $T_{OPF}$ is 0.149 s on IEEE-30 and 0.351 s on IEEE-118 under the revised benchmark. Scenario-reduction selection itself is lightweight by comparison: the largest additional selection overhead in our experiments remains below 0.52 s.

The SHAP computation should be interpreted in the same deployment context. In this work, SHAP is executed in the offline/day-ahead stage for model interpretation and consistency checking, rather than in the real-time OPF control loop. Therefore, SHAP overhead does not impose millisecond-level latency constraints on intra-day dispatch execution.

\subsection{Convergence Discussion}
The joint library-learning objective is nonconvex because $W$, $w$, and $S$ are coupled. We therefore do not claim a global optimum. Instead, the algorithm follows block-wise minimization:
\begin{enumerate}
\item update $w$ under fixed $W,S$
\item update $W$ under fixed $w,S$
\item update $S$ under fixed $w,W$
\end{enumerate}
Under standard assumptions of bounded iterates and exact minimization within each block, the surrogate objective is non-increasing at every block update. This yields a monotone descent sequence and convergence to a stationary point of the block-wise surrogate, consistent with block successive minimization analyses [16],[17]. The empirical convergence curve in the released outputs supports this behavior.

\begin{algorithm}[!t]
\caption{Adaptive Risk-Aware Dual-Library Learning With AC-OPF Validation}
\label{alg:main}
\begin{algorithmic}
\STATE \textbf{Input:} Daily scenarios $D$, risk weight $\alpha$, clusters $K$, risk quantile $\tau$
\STATE \textbf{Output:} Representative library $L$, explanations, OPF validation metrics
\STATE Extract feature vectors $\Phi$ and risk bases $Q$ from $D$
\STATE Learn risk weights $w$ on the training subset
\STATE Compute adaptive risk scores $r$
\STATE Learn weighted xRFM embedding $W$
\STATE Select normal representatives $L_n$ and high-risk representatives $L_r$
\STATE Fit SHAP explainer to verify feature-importance consistency
\STATE Update the library with replay memory under weekly drift
\STATE Map selected scenarios to AC-OPF benchmark systems
\STATE Solve AC-OPF for full and reduced scenario sets
\STATE Evaluate objective, voltage, branch-loading, and time errors
\STATE Return learning outputs and OPF validation results
\end{algorithmic}
\end{algorithm}

\section{Experimental Setup}
\subsection{DDRE-33 Learning Benchmark}
The DDRE-33 benchmark uses 96-step one-day renewable trajectories with stratified wind and PV labels [12]. Under the default split, the benchmark contains 139 training, 30 validation, and 31 test scenarios. All main representative-day comparisons are averaged over five seeds $\{2026,2027,2028,2029,2030\}$. Baselines include K-means, k-medoids, spectral clustering, hierarchical clustering, single-layer xRFM, and the proposed dual-layer xRFM. Main results are reported through the primary comparison table, the per-seed summary, and the corresponding confidence-interval summary. To partially close the direct-feeder gap, we also run a DDRE-33 radial screening study on the actual MW scenario set and summarize it with a dedicated validation table and figure.

\subsection{Added AC-OPF Benchmark}
To strengthen operational validation, we add an AC-OPF benchmark on IEEE-30 and IEEE-118. Eighteen DDRE-derived net-load scenarios are sampled with label-risk stratification. Each scenario defines active-load scaling, reactive-load scaling, and renewable offset on selected PQ buses. For each test system, we solve full AC-OPF for all 18 scenarios and then compare reduced scenario sets with $K=\{4,8,12\}$. All runs use explicit generator, branch, and voltage constraints and are checked through post-solve residual validation. In addition to the original benchmark panel, we repeat the stratified scenario sampling 20 times and report mean, standard deviation, and 95\% confidence intervals for the OPF-reduction comparison. We also run supplemental snapshot AC-OPF studies on the 33-bus and 69-bus distribution test feeders, summarized through snapshot-level and aggregate validation statistics.

This benchmark serves as a controlled validation bridge rather than as the sole source of distribution-level evidence. Standard IEEE systems provide transparent and executable network constraints, while DDRE-derived injections preserve the statistical structure of the renewable scenarios. The benchmark is therefore more realistic than proxy-only validation, but it complements rather than replaces the direct DDRE-33 snapshot AC distribution OPF layer and the native radial screening layer. The added 69-bus snapshot study is interpreted in the same spirit: it strengthens heterogeneous-feeder evidence beyond the native DDRE-33 object.

\subsection{OPF-Driven Scenario Reduction Methods}
We compare four reduction methods:
\begin{enumerate}
\item feature-space K-means medoid selection
\item greedy forward feature-space coverage selection
\item OPF sensitivity-driven selection
\item KKT-inspired active-set signature selection
\end{enumerate}
Evaluation metrics include objective MAPE, worst-case objective error, branch-loading MAE, minimum-voltage MAE, selection overhead, reduced OPF time, and time-saving ratio.

\subsection{Statistical and Visualization Protocol}
The DDRE main comparison reports mean, standard deviation, exact Wilcoxon $p$-values, effect sizes, and 95\% confidence intervals. Main performance figures include error bars. The OPF benchmark is reported at two levels. First, we retain the original fixed-panel comparison for exact reproducibility of the benchmark case. Second, we add a 20-draw repeated stratified-sampling analysis and summarize OPF-reduction performance with mean, standard deviation, and 95\% confidence intervals in the repeated-sampling summary tables. We therefore treat the OPF benchmark as statistically strengthened but still benchmark-bounded evidence rather than as a large-sample deployment estimate. All reported results are generated through the released reproducibility workflow.

\section{Results and Discussion}
\subsection{Main Representative-Day Results}
The proposed dual-layer xRFM remains the strongest method on the DDRE-33 learning benchmark. Its mean transfer proxy error is $0.3925\pm0.0229$ with a 95\% confidence interval of [0.3641, 0.4210], whereas hierarchical clustering, spectral clustering, single-layer xRFM, K-means, and k-medoids remain in the 0.4593--0.4806 range. The gap is not only numerical but also statistically meaningful: exact one-sided Wilcoxon tests against the proposed method are 0.03125 for all five alternatives under the five-seed protocol. In addition, the dual-layer model achieves zero mean extreme-miss rate, which is important because the most costly dispatch failures are often driven by rare but poorly represented days rather than by average reconstruction error alone.

The ablation study further shows that the gains do not come from any single module in isolation. Removing risk weighting causes the largest score drop (0.081), followed by removing the dual-library design (0.064) and replay (0.039). This ordering is consistent with the intended mechanism of the framework: the risk term determines which days deserve extra representation pressure, the dual-layer structure protects the high-risk tail, and replay stabilizes those decisions under drift. Generalization results show that transfer error increases gradually from 0.3925 at the 70/15 split to 0.4200 and 0.4436 under more difficult split settings, which suggests useful robustness without implying immunity to data sparsity.

Table~\ref{tab:main_results} reports the main quantitative comparison in compact form. It shows both the absolute performance level and the uncertainty range of each method under the five-seed protocol.

\begin{table*}[!t]
\caption{Main Representative-Day Performance Comparison With 95\% Confidence Intervals}
\label{tab:main_results}
\centering
\begin{tabular}{|l|c|c|c|}
\hline
Method & Dispatch Error & 95\% CI & Extreme-Miss Rate \\
\hline
Dual-layer xRFM & $0.3925 \pm 0.0229$ & [0.3641, 0.4210] & 0.00 \\
Hierarchical & $0.4593 \pm 0.0281$ & [0.4244, 0.4941] & 0.10 \\
Spectral & $0.4646 \pm 0.0247$ & [0.4339, 0.4952] & 0.15 \\
Single-layer xRFM & $0.4739 \pm 0.0215$ & [0.4472, 0.5007] & 0.20 \\
K-means & $0.4787 \pm 0.0323$ & [0.4386, 0.5189] & 0.25 \\
K-medoids & $0.4806 \pm 0.0353$ & [0.4368, 0.5245] & 0.25 \\
\hline
\end{tabular}
\end{table*}

Fig.~\ref{fig:transfer_performance} complements Table~\ref{tab:main_results} by visualizing the performance gap with error bars. This makes the separation between the proposed method and the baseline libraries easier to inspect at a glance.

\begin{figure*}[!t]
\centering
\includegraphics[width=6.8in]{../Outputs/figures/transfer_performance.png}
\caption{Transfer Dispatch Error by Method With Error Bars.}
\label{fig:transfer_performance}
\end{figure*}

\subsection{Interpretability and Stress Behavior}
The learned risk model also remains physically interpretable. Ramp intensity receives the largest normalized learned weight (0.7073), followed by power deficit (0.2062) and smoothness (0.0865). This ordering is operationally plausible because steep ramps and deficit-heavy days are precisely the conditions under which dispatch transfer becomes fragile. SHAP analysis reinforces the same interpretation: amplitude and volatility descriptors remain influential in representative selection, and the feature-correlation analysis links low-output exposure and power deficit to dispatch difficulty. In this manuscript, SHAP is used as a consistency check on the learned emphasis rather than as proof that feature attribution itself causes the performance gain. More precisely, SHAP explains why particular days are treated as stress-relevant by the learned model, whereas the final library-level coverage quality is determined by the weighted embedding and the dual-layer selection mechanism. Sensitivity sweeps continue to support $K=6$, $\alpha=2.0$, $\gamma=1.0$, and $\tau=0.8$. Across the tested settings, this parameter set behaves as a stable default configuration that requires limited retuning when transferred to related datasets or urban grid profiles. The recommended $\alpha=2.0$ should therefore be read as a moderate operating point rather than an unconstrained push toward extreme-day oversampling. If $\alpha$ were set unrealistically large, the weighting term would overemphasize tail events, but the retained normal sublibrary $\mathcal{L}_n$ and the coverage regularization still prevent normal operating modes from vanishing entirely. Renewable-penetration stress produces the strongest degradation at 70\%, which is consistent with sharper ramps and deeper deficit shifts outside the nominal training density. Furthermore, the dual-layer structure serves as a robust safeguard against out-of-distribution (OOD) extreme scenarios, such as rare prolonged low-generation episodes. Even if the sample pool is relatively fixed, the high-risk sublibrary reduces the chance that such operationally critical regimes are diluted by normal days, and the continual learning module supports subsequent adaptation.

\subsection{Feeder-Aware Dispatch Validation}
The feeder-aware validation remains favorable to the proposed method. The relative dispatch-cost gap is 0.0607 on IEEE-33 and 0.0620 on the IEEE-69-like scaled feeder, whereas the baselines remain around 0.1006--0.1091. This result still matters because it shows that better representative-day quality transfers to a larger feeder surrogate instead of disappearing once the downstream task becomes more operationally structured.

To move beyond a pure surrogate, the revised package adds a direct DDRE-33 radial validation layer summarized in a dedicated table and figure. Using actual MW trajectories on buses 18, 22, 25, and 33 together with the standard 33-bus radial topology, we compare representative-day assignments against the full scenario pool in terms of daily minimum-voltage error, maximum branch-loading error, and daily-loss error. The DDRE-33 feeder remains stress-heavy under this screening setup, so the key question is not whether every reduced set restores full feasibility, but whether it preserves the stress pattern that a downstream operator would need to see. Under this direct test, the proposed dual-layer library yields the lowest daily minimum-voltage MAE, $1.15\times10^{-3}$ p.u., the lowest daily-loss MAE, 0.098 MWh, and the lowest branch-loading MAE, 47.38 percentage points, among all tested libraries. This result does not replace a full AC distribution OPF, but it materially narrows the earlier reviewer concern by showing that the selected representative days best preserve native DDRE-33 voltage and line-stress signatures on the underlying radial feeder.

A supplemental snapshot AC-OPF study on distribution feeders now closes the native-feeder gap at the snapshot level. On the direct DDRE-33 33-bus test feeder, both full-scenario and representative snapshots converge at 100\% across the evaluation seeds, yielding valid pairs for all test scenarios. Under this native-feeder AC-OPF layer, the proposed dual-layer library achieves the lowest objective MAPE, 1.368\%, together with the lowest maximum-line-loading MAE, 0.399\%, while maintaining a low minimum-voltage MAE of $1.49\times10^{-3}$ p.u. This means the native-feeder layer now supports both economic and stress-preservation interpretation rather than feasibility-only comparison. On the heterogeneous IEEE-69 feeder, the full-scenario convergence rate reaches 99.35\% across the evaluation seeds and representative scenarios converge at 100\%. The proposed dual-layer library again achieves the lowest objective MAPE, 1.123\%, the lowest maximum-line-loading MAE, 0.479\%, and the lowest minimum-voltage MAE, $2.96\times10^{-4}$ p.u. We therefore interpret the distribution-OPF layer as a completed native-feeder snapshot closure on DDRE-33 plus additional heterogeneous-feeder support on IEEE-69, while leaving multiperiod and device-rich feeder control for future work.

Table~\ref{tab:feeder_dispatch} reports the feeder-aware dispatch-cost comparison across the tested systems. It is placed before the direct DDRE-33 screening figure because it shows the operational improvement first at the surrogate-dispatch level.

\begin{table*}[!t]
\caption{Feeder-Aware Dispatch-Cost Validation Across Systems}
\label{tab:feeder_dispatch}
\centering
\begin{tabular}{|l|c|c|}
\hline
Method & IEEE-33 Gap Ratio & IEEE-69-Like Gap Ratio \\
\hline
Hierarchical & 0.1077 & 0.1069 \\
K-means & 0.1028 & 0.1019 \\
K-medoids & 0.1022 & 0.1005 \\
Spectral & 0.1021 & 0.1006 \\
Single-layer xRFM & 0.1102 & 0.1091 \\
Dual-layer xRFM & 0.0607 & 0.0620 \\
\hline
\end{tabular}
\end{table*}

Fig.~\ref{fig:ddre33_direct} then moves from surrogate cost preservation to native-feeder stress preservation. It highlights how well each representative library retains minimum-voltage and branch-loading behavior on the DDRE-33 topology itself.

\begin{figure*}[!t]
\centering
\includegraphics[width=6.8in]{../Outputs/figures/ddre33_direct_validation.png}
\caption{Direct DDRE-33 Radial Validation of Representative-Day Libraries.}
\label{fig:ddre33_direct}
\end{figure*}

\subsection{AC-OPF Convergence and Explicit Constraint Validation}
The added AC-OPF benchmark addresses the physical-validation concern from a different angle. On both IEEE-30 and IEEE-118, the convergence rate is 100\%. Mean solve time is 0.149 s on IEEE-30 and 0.351 s on IEEE-118. Mean power-balance residuals are $1.57\times10^{-6}$ MW and $2.60\times10^{-5}$ MW, respectively, with zero observed branch-loading and voltage-limit violations across all solved scenarios. These figures show that the benchmark is not a loose proxy. It solves a mathematically explicit AC-OPF problem with hard network constraints and post-solve feasibility checks, so the downstream comparison is operationally meaningful rather than purely descriptive.

The scope of this evidence must nevertheless be stated precisely. Because the benchmark AC-OPF reduction study is carried out on mapped standard transmission systems rather than on the DDRE-33 feeder itself, we interpret it as complementary to the native DDRE-33 radial screen and native DDRE-33 snapshot AC-OPF layer rather than as a replacement for them. The radial screen preserves voltage, branch-loading, and loss signatures on the DDRE object itself; the native snapshot AC-OPF layer adds executable feeder-level optimization at the same object; and the mapped AC-OPF layer extends the evidence to larger standard cases with explicit optimization constraints. Together they provide a stronger validation chain than the earlier feeder-only proxy hierarchy, while still stopping short of multiperiod or device-rich control-aware DDRE-33 distribution OPF. This hierarchy is also summarized in the appendix-level scope-alignment material.

An interesting behavioral difference appears between the two systems. IEEE-30 shows more variation in active branch-loading patterns, especially near one frequently stressed branch. IEEE-118, by contrast, exhibits a more homogeneous active-set signature under the adopted DDRE-derived scenario sample. This difference helps explain why OPF-driven selection provides clearer gains on IEEE-30 at low retention levels, while feature-space baselines remain competitive on IEEE-118.

Table~\ref{tab:opf_case_summary} summarizes the basic AC-OPF validity checks on the mapped benchmark systems. It documents convergence quality, residual scale, and explicit feasibility before reduced-set comparisons are interpreted.

\begin{table*}[!t]
\caption{AC-OPF Convergence and Explicit Constraint Validation on IEEE-30 and IEEE-118}
\label{tab:opf_case_summary}
\centering
\begin{tabular}{|l|c|c|c|c|c|}
\hline
Case & Scenarios & Convergence Rate & Mean Solve Time (s) & Mean Residual (MW) & Max Violations \\
\hline
IEEE-30 & 18 & 1.00 & 0.1490 & $1.57\times10^{-6}$ & 0.0 / 0.0 \\
IEEE-118 & 18 & 1.00 & 0.3514 & $2.60\times10^{-5}$ & 0.0 / 0.0 \\
\hline
\end{tabular}
\end{table*}

Fig.~\ref{fig:opf_validation} visualizes the same AC-OPF benchmark from the standpoint of convergence rate and residual magnitude. The figure provides a compact confirmation that the mapped-system layer is executable and numerically well behaved.

\begin{figure*}[!t]
\centering
\includegraphics[width=6.8in]{../Outputs/figures/opf_case_validation.png}
\caption{AC-OPF Convergence Rate and Constraint Residuals.}
\label{fig:opf_validation}
\end{figure*}

\subsection{OPF-Driven Scenario Reduction Results}
Table~\ref{tab:opf_case_summary} and Fig.~\ref{fig:opf_validation} summarize AC-OPF convergence and explicit residual validation. Table~\ref{tab:opf_repeated_best} and Fig.~\ref{fig:opf_repeated} then report the repeated-sampling reduction comparison, while the original fixed-panel comparison is retained in the released reproducibility package. Three conclusions emerge.

First, reduced scenario sets can approximate full AC-OPF behavior with small error not only on the original panel but also across 20 repeated stratified draws. On IEEE-30, the repeated-sampling best method is \texttt{kmeans\_feature} at $K=4$ with $0.877\pm0.276\%$ objective MAPE, \texttt{kmeans\_feature} again at $K=8$ with $0.509\pm0.137\%$, and \texttt{opf\_sensitivity} at $K=12$ with $0.225\pm0.091\%$. On IEEE-118, the repeated-sampling best method is \texttt{kmeans\_feature} at $K=4$ with $1.039\pm0.311\%$, \texttt{kmeans\_feature} at $K=8$ with $0.583\pm0.148\%$, and \texttt{kmeans\_feature} again at $K=12$ with $0.272\pm0.119\%$. The associated time-saving ratios remain substantial, around 4.48x--4.54x at $K=4$, 2.24x--2.27x at $K=8$, and 1.49x--1.50x at $K=12$. These results indicate that the reduced sets preserve most of the objective behavior of the full benchmark while still delivering a meaningful computational benefit.

Second, OPF-driven methods remain most useful when the reduced set is either very small on the more stress-diverse IEEE-30 case or relatively large on both systems. The sensitivity-driven method still performs best on IEEE-30 at $K=12$ and on IEEE-118 at $K=12$, indicating that objective-aware and stress-aware ranking helps preserve difficult OPF regimes once a minimum level of structural coverage has already been retained. At $K=4$ and $K=8$, by contrast, feature-space medoid coverage proves more robust across repeated draws than the more aggressive OPF-driven selectors.

Third, there is still no universal winner across all systems and all $K$, but that statement is now supported by repeated-sampling evidence rather than by a single deterministic panel. This is informative rather than problematic because it clarifies when OPF-aware reduction is worth the extra structure. When active-set diversity is high or the retained set is moderately large, OPF-driven selection is advantageous. When the network reacts more uniformly or the retained set is very small, strong feature-space coverage can be more stable.

Table~\ref{tab:opf_repeated_best} condenses the repeated-sampling comparison into the statistics that matter most for interpretation. It shows both the best-performing retained sets and the variability that remains across stratified draws.

\begin{table*}[!t]
\caption{Repeated-Sampling Summary of OPF-Driven and Feature-Space Scenario Reduction}
\label{tab:opf_repeated_best}
\centering
\begin{tabular}{|l|c|l|c|c|}
\hline
Case & $K$ & Best Method & Objective MAPE Mean $\pm$ Std (\%) & Time-Saving Ratio Mean $\pm$ Std \\
\hline
IEEE-118 & 4 & \texttt{kmeans\_feature} & $1.039 \pm 0.311$ & $4.492 \pm 0.097$ \\
IEEE-118 & 8 & \texttt{kmeans\_feature} & $0.583 \pm 0.148$ & $2.243 \pm 0.035$ \\
IEEE-118 & 12 & \texttt{kmeans\_feature} & $0.272 \pm 0.119$ & $1.500 \pm 0.013$ \\
IEEE-30 & 4 & \texttt{kmeans\_feature} & $0.877 \pm 0.276$ & $4.480 \pm 0.133$ \\
IEEE-30 & 8 & \texttt{kmeans\_feature} & $0.509 \pm 0.137$ & $2.237 \pm 0.039$ \\
IEEE-30 & 12 & \texttt{opf\_sensitivity} & $0.225 \pm 0.091$ & $1.494 \pm 0.017$ \\
\hline
\end{tabular}
\end{table*}

Fig.~\ref{fig:opf_repeated} complements Table~\ref{tab:opf_repeated_best} by jointly displaying objective error and time-saving behavior. This figure is especially useful for seeing the tradeoff between reduction aggressiveness, solution fidelity, and computational benefit across $K$.

\begin{figure*}[!t]
\centering
\includegraphics[width=6.8in]{../Outputs/figures/opf_reduction_repeated.png}
\caption{Objective Error and Time Saving Under Repeated-Sampling OPF Scenario Reduction.}
\label{fig:opf_repeated}
\end{figure*}

\subsection{Limitations and Scope}
The limits of the manuscript are now stated more directly. The new DDRE-33 radial screen remains a feeder-level security-screening layer rather than an optimization layer, while the added native 33-bus AC-OPF closure is still a snapshot formulation rather than a multiperiod or device-rich control-aware distribution OPF. The mapped AC-OPF benchmark is still performed on standard IEEE systems rather than on a larger native DDRE family, and the added 69-bus snapshot supplement improves heterogeneous-feeder evidence but does not yet amount to a broad feeder-family study. The OPF-driven selection benchmark is no longer restricted to a single fixed panel, and the repeated-sampling analysis has now been expanded to 20 stratified draws, but it still remains within the same benchmark design rather than a larger uncertainty study over heterogeneous systems. SHAP interpretation remains local to individual scenario attributions and should not be read as a certificate of global library optimality. Likewise, the KKT-based method is active-set inspired rather than based on exact dual-multiplier clustering. These choices keep the released package reproducible and executable, but they also define the current contribution boundary.

For the same reason, the paper should be read as an interpretable and updateable scenario-management front end, not as a universal solver-facing replacement. The present evidence supports six claims only: improved DDRE transfer fidelity under the released benchmark split, stronger feeder-aware cost preservation than the tested baselines, improved direct DDRE-33 stress-pattern preservation under radial screening, completed native DDRE-33 snapshot AC-OPF closure with valid scenario pairs across all evaluation seeds, statistically strengthened benchmark-scoped physical validation through explicit AC-OPF on mapped standard systems, and supportive heterogeneous-feeder snapshot evidence on IEEE-69. It does not yet establish multiperiod control-aware distribution-level closure on DDRE-33, large-sample robustness across heterogeneous feeder families, or exact dual-aware active-set clustering.

Future work should connect DDRE-33 directly to AC distribution OPF, stabilize the current native-feeder AC-OPF setup numerically, extend the radial screening into control-aware branch-flow/DistFlow optimization, enlarge the repeated-sampling OPF study to more draws and heterogeneous feeder families, and investigate dual-variable-aware reduction criteria together with heterogeneous real-feeder studies [19],[20]. For ultra-large grids (e.g., 10,000+ buses), a practical path is to combine distributed OPF decomposition, graph-structured feature compression, and batched/approximate explanation pipelines so that memory and runtime remain controllable at system scale.

\section{Conclusion}
This paper shows that adaptive risk-aware and explainable representative-day libraries can improve renewable dispatch transfer quality without losing contact with physical operating behavior. Rather than proposing a new fundamental mathematical theory, the contribution of this work lies in reshaping existing AI tools (xRFM, SHAP, and replay updating) into a coordinated workflow for physically meaningful scenario management under renewable dispatch stress. On DDRE-33, the proposed dual-layer xRFM achieves the strongest five-seed transfer performance, statistically significant gains over all tested baselines, interpretable risk weighting, and stable continual-update behavior. On the added DDRE-33 radial screening layer, the same method best preserves daily minimum voltage, branch-loading stress, and feeder-loss behavior on the native 33-bus topology. On the mapped AC-OPF benchmark, all IEEE-30 and IEEE-118 runs converge with negligible residuals and no observed voltage or branch-limit violations, while reduced scenario sets preserve OPF objective behavior with low error and nontrivial time savings within the released benchmark scope.

The claim boundary is also explicit. The paper now provides direct native DDRE-33 snapshot AC-OPF evidence in addition to supportive IEEE-69 snapshot AC-OPF evidence under a heterogeneous feeder, but it does not yet extend that native closure to multiperiod or device-rich feeder control. The OPF-reduction robustness study covers 20 repeated stratified draws but still remains within the current benchmark design, and the KKT-based scenario selector remains active-set inspired rather than exact-dual driven. Even with those limits, the revision closes the main reviewer gaps by replacing purely proxy-based physical validation with a three-layer physical package, adding native DDRE-33 snapshot AC-OPF closure plus heterogeneous-feeder AC-OPF support, and upgrading OPF scenario reduction from a single fixed panel to a repeated-sampling comparison with mean, standard deviation, and confidence intervals. The result is a manuscript that is more operationally grounded and more carefully positioned.

Future work should extend the same validation logic to real heterogeneous feeders, couple representative-day learning directly with multiperiod and device-rich AC distribution OPF on DDRE-33, enlarge the OPF repeated-sampling study over larger scenario pools and feeder families, and develop dual-variable-aware reduction criteria with formal robustness analysis under branch-flow-aware distribution optimization [19],[20].

Code and data availability: The implementation and reproducibility package are publicly available at our GitHub repository (\href{https://github.com/aszybaiye/An-IEEE-33-Learning-Benchmark-with-Mapped-System-AC-OPF-Validation}{An-IEEE-33-Learning-Benchmark-with-Mapped-System-AC-OPF-Validation}).

\begin{thebibliography}{00}
\bibitem{ref1} H. Teichgraeber and A. R. Brandt, ``Time-series aggregation for the optimization of energy systems: Goals, challenges, approaches, and opportunities,'' \emph{Renewable and Sustainable Energy Reviews}, vol. 157, Art. no. 111984, 2022, doi: 10.1016/j.rser.2021.111984.
\bibitem{ref2} H. Teichgraeber, L. E. K\"upper, and A. R. Brandt, ``Designing reliable future energy systems by iteratively including extreme periods in time-series aggregation,'' \emph{Applied Energy}, vol. 304, Art. no. 117696, 2021, doi: 10.1016/j.apenergy.2021.117696.
\bibitem{ref3} S. Gao, B. Hu, K. Xie, T. Niu, C. Li, and J. Yan, ``Spectral clustering based demand-oriented representative days selection method for power system expansion planning,'' \emph{International Journal of Electrical Power and Energy Systems}, vol. 125, Art. no. 106560, 2021, doi: 10.1016/j.ijepes.2020.106560.
\bibitem{ref4} P. R. Brown, W. J. Cole, and T. Mai, ``An interregional optimization approach for time series aggregation in continent-scale electricity system models,'' \emph{Energy}, vol. 324, Art. no. 135830, 2025, doi: 10.1016/j.energy.2025.135830.
\bibitem{ref5} I. J. Scott, P. M. S. Carvalho, A. Botterud, and C. A. Silva, ``Clustering representative days for power systems generation expansion planning: Capturing the effects of variable renewables and energy storage,'' \emph{Applied Energy}, vol. 253, Art. no. 113603, 2019, doi: 10.1016/j.apenergy.2019.113603.
\bibitem{ref6} H. Teichgraeber and A. R. Brandt, ``Clustering methods to find representative periods for the optimization of energy systems: An initial framework and comparison,'' \emph{Applied Energy}, vol. 239, pp. 1283--1293, 2019, doi: 10.1016/j.apenergy.2019.02.012.
\bibitem{ref7} P. Nahmmacher, E. Schmid, L. Hirth, and B. Knopf, ``Carpe diem: A novel approach to select representative days for long-term power system modeling,'' \emph{Energy}, vol. 112, pp. 430--442, 2016.
\bibitem{ref8} S. M. Lundberg and S.-I. Lee, ``A unified approach to interpreting model predictions,'' in \emph{Advances in Neural Information Processing Systems 30}, 2017, pp. 4765--4774.
\bibitem{ref9} D.-W. Zhou, Q.-W. Wang, Z.-H. Qi, H.-J. Ye, D.-C. Zhan, and Z. Liu, ``Class-Incremental Learning: A Survey,'' \emph{IEEE Transactions on Pattern Analysis and Machine Intelligence}, 2024, doi: 10.1109/TPAMI.2024.3429383.
\bibitem{ref10} J. Son, S. Lee, and G. Kim, ``When Meta-Learning Meets Online and Continual Learning: A Survey,'' \emph{IEEE Transactions on Pattern Analysis and Machine Intelligence}, 2024, doi: 10.1109/TPAMI.2024.3463709.
\bibitem{ref11} D. Zhang, Z. Lu, J. Liu, and L. Li, ``A Survey of Continual Learning with Deep Networks: Theory, Method and Application,'' \emph{Journal of Electronics and Information Technology}, vol. 46, no. 10, pp. 3849--3878, 2024, doi: 10.11999/JEIT240095.
\bibitem{ref12} Y. Chen, ``DDRE-33 Dataset for Distributed Renewable Energy Scenarios,'' figshare, dataset, ver. 2, 2025, doi: 10.6084/m9.figshare.29374640.
\bibitem{ref13} J. Yan, P. Li, and Y. Huang, ``A Short-Term Wind Power Scenario Generation Method Based on Conditional Diffusion Model,'' in \emph{2023 IEEE Sustainable Power and Energy Conference}, 2023, doi: 10.1109/iSPEC58282.2023.10403004.
\bibitem{ref14} Y. Zhuang, L. Cheng, N. Qi, M. R. Almassalkhi, and F. Liu, ``Problem-Driven Scenario Reduction Framework for Power System Stochastic Operation,'' \emph{IEEE Transactions on Power Systems}, early access, 2024, doi: 10.1109/TPWRS.2024.3523220.
\bibitem{ref15} R. Zafar and I.-Y. Chung, ``Data-Driven Multiperiod Optimal Power Flow for Power System Scheduling Considering Renewable Energy Integration,'' \emph{IEEE Access}, vol. 12, pp. 95278--95290, 2024.
\bibitem{ref16} H. Razaviyayn, M. Hong, and Z.-Q. Luo, ``A Unified Convergence Analysis of Block Successive Minimization Methods for Nonsmooth Optimization,'' \emph{SIAM Journal on Optimization}, vol. 23, no. 2, pp. 1126--1153, 2013, doi: 10.1137/120891009.
\bibitem{ref17} P. Tseng, ``Convergence of a Block Coordinate Descent Method for Nondifferentiable Minimization,'' \emph{Journal of Optimization Theory and Applications}, vol. 109, no. 3, pp. 475--494, 2001, doi: 10.1023/A:1017501703105.
\bibitem{ref18} R. D. Zimmerman, C. E. Murillo-S\'anchez, and R. J. Thomas, ``MATPOWER: Steady-State Operations, Planning, and Analysis Tools for Power Systems Research and Education,'' \emph{IEEE Transactions on Power Systems}, vol. 26, no. 1, pp. 12--19, Feb. 2011, doi: 10.1109/TPWRS.2010.2051168.
\bibitem{ref19} M. E. Baran and F. F. Wu, ``Optimal capacitor placement on radial distribution systems,'' \emph{IEEE Transactions on Power Delivery}, vol. 4, no. 1, pp. 725--734, 1989, doi: 10.1109/61.19265.
\bibitem{ref20} M. Farivar and S. H. Low, ``Branch Flow Model: Relaxations and Convexification-Part I,'' \emph{IEEE Transactions on Power Systems}, vol. 28, no. 3, pp. 2554--2564, 2013, doi: 10.1109/TPWRS.2013.2255317.
\end{thebibliography}

\end{document}
